\Needspace{16\baselineskip}
\fancyhead[L]{\small VESTA. JUNO. PALLAS. CERES.}
\begin{center}
VESTA. JUNO. PALLAS. CERES.
\end{center}

\begin{center}
\rule{4cm}{0.4pt}\\
\fontsize{10}{12}\selectfont Göttingische gelehrte Anzeigen. Stück 107. Seite 1065..1068. 4. Juli 1808.\\
\rule{4cm}{0.4pt}
\end{center}

In der Nacht vom 22. auf den 23. Juni wurde der neue Planet \emph{Vesta}, nachdem er bei uns fast ein ganzes Jahr unsichtbar gewesen war, an der hiesigen Sternwarte wiedergefunden und zum ersten Mal wieder beobachtet. Er gleicht bereits ganz einem Stern achter Größe, und während der Beobachtungen, die eine Stunde bis zum Tagesanbruch fortgesetzt wurden, war die Änderung seiner Position schon deutlich zu erkennen. Das Mittel aus allen Vergleichen mit einem Stern achter Größe, dessen scheinbare Position aus der Hist. céleste reduziert und mit


\begin{center}\fontsize{10}{12}\selectfont\setlength{\tabcolsep}{3pt}\renewcommand{\arraystretch}{1.15}\begin{tabular}{l}
$352^\circ\;44'\;1''$ Rektaszension,\\
$8^\circ\;58'\;9''$ südliche Deklination
\end{tabular}\end{center}


\noindent angesetzt wurde, ergab folgendes Resultat:


\begin{center}\fontsize{10}{12}\selectfont\setlength{\tabcolsep}{3pt}\renewcommand{\arraystretch}{1.15}\begin{tabular}{lrrr}
1808 & \shortstack{Mittlere Zeit\\in\\Göttingen} & \shortstack{Scheinbare\\Rektaszension\\der Vesta} & \shortstack{Südliche\\Deklination\\der Vesta}\\\hline
Juni 22 & $13^{\text{u}}\,46^{\text{m}}\,25^{\text{s}}$ & $353^\circ\;12'\;17''$ & $9^\circ\;11'\;55''$\\
\end{tabular}\end{center}


\par\indent Die nach den dritten Elementen von Professor \textsc{Gauss} berechnete und im Septemberheft der Monatl. Correspondenz 1807 abgedruckte Ephemeride gibt die Rektaszension um $8'$ zu klein und die Deklination um $3'$ zu groß an. Diese Elemente waren schon im Juli des Vorjahrs aus allen damals verfügbaren Beobachtungen berechnet worden. Die letzte war am 11. Juli in Lilienthal gemacht worden. Die Mailänder Astronomen verfolgten den Planeten noch viel länger (bis zum 26. Sept.). Ihre Beobachtungen hätten im Nachhinein eine noch viel genauere Bestimmung der Bahn ermöglicht. Obwohl diese letzten Beobachtungen nach der Berechnung von Professor \textsc{Harding} bereits um mehr als eine Minute von den dritten Elementen abwichen, hielt Professor \textsc{Gauss} damals eine erneute Korrektur der Elemente für unnötig. Die schon berechnete Ephemeride hielt er für mehr als genau genug, um den Planeten wiederzufinden, wie der Erfolg nun auch bestätigt hat. Später fand Professor \textsc{Gauss} doch noch einen Anlass für eine neue Berechnung der Elemente. Er wollte einen Abschnitt seines im Druck befindlichen Werks über die Bahnen der Himmelskörper durch ein ausgewähltes Beispiel erläutern. Das genauestmögliche Endergebnis der Beobachtungen des Vorjahrs lag außerhalb dieses Ziels. Dabei wurde eine der letzten Mailändi\SourceLeaf{309}schen Beobachtungen mitverwendet, dem Zweck der Rechnung entsprechend jedoch \emph{ohne} sie vorher vom wahrscheinlichen Beobachtungsfehler zu befreien. Trotzdem war abzusehen, dass das Ergebnis beträchtlich genauer sein würde als die dritten Elemente. Da es bisher noch nirgends gedruckt ist, fügen wir es hier bei:

\begin{center}
\emph{Elemente der Vesta.}
\end{center}


\begin{center}\fontsize{9}{11}\selectfont\setlength{\tabcolsep}{3pt}\renewcommand{\arraystretch}{1.15}\begin{tabular}{lr}
Epoche 1807 Meridian von Paris .\;.\;.\;.\;.\;.\;.\;.\; & $168^\circ\;10'\;47''6$\\
Tägliche tropische mittlere Bewegung .\;.\;.\;.\;.\;.\;.\;.\; & $978''8588$\\
Perihel .\;.\;.\;.\;.\;.\;.\;.\;.\;.\;.\;.\;.\;.\;.\; & $249^\circ\;57'\;6''5$\\
Aufsteigender Knoten .\;.\;.\;.\;.\;.\;.\;.\;.\;,\;.\;.\; & $103\;11\;57$\\
\multicolumn{2}{l}{\quad beide für 1807 und gegenüber den Fixsternen unverändert}\\
Neigung der Bahn .\;.\;.\;.\;.\;.\;.\;.\;.\;.\;.\;.\;.\; & $7\;8\;21$\\
Exzentrizität .\;.\;.\;.\;.\;.\;.\;.\;.\;.\;,\;.\;.\;.\;.\;.\; & $0{,}0880158$\\
Logarithmus des mittleren Abstands .\;.\;.\;\hspace{.5em}.\;.\;.\;.\; & $0{,}3728980$\\
\end{tabular}\end{center}


\Needspace{9\baselineskip}
\par\indent Die mit \emph{diesen} Elementen für den Zeitpunkt der obigen Beobachtung berechnete Position lautet:


\begin{center}\fontsize{10}{12}\selectfont\setlength{\tabcolsep}{3pt}\renewcommand{\arraystretch}{1.15}\begin{tabular}{rrrr}
Rektaszension & Unterschied & Deklination & Unterschied\\\hline
$353^\circ\;14'\;30''$ & $+2'\;13''$ & $9^\circ\;12'\;26''$ & $+31''$\\
\end{tabular}\end{center}


\par\indent Es ist noch zu bemerken, dass der Vergleichsstern für die Bestimmung der Deklination sehr ungünstig stand und diese daher nicht sehr zuverlässig ist. Für die Rektaszension glauben wir umso mehr einstehen zu können, da die Hist. cél. zwei Beobachtungen des Vergleichssterns enthält, die sehr gut übereinstimmende Ergebnisse lieferten. Professor \textsc{Gauss} hat die Refraktion bei der Reduktion nach einem eigenen Verfahren berücksichtigt.

\par\indent Schon am 20. Juni wurde hier auch \emph{Juno} wiedergefunden. Sie ist noch sehr lichtschwach und gleicht nur einem Stern 10..11. Größe. Folgende Beobachtungen waren bisher möglich: [Sie stehen oben auf Seite 302 im Brief an \textsc{von Lindenau} vom 27. Juni 1808.]

\par\indent Die Position des Vergleichssterns wurde aus der Histoire céleste bestimmt:


\begin{center}\fontsize{10}{12}\selectfont\setlength{\tabcolsep}{3pt}\renewcommand{\arraystretch}{1.15}\begin{tabular}{l}
Rektaszension $315^\circ\;24'\;58''$\\
südliche Deklination $2^\circ\;4'\;20''$
\end{tabular}\end{center}


\par\indent Demnach weichen die siebten, im 14. Stück dieser Anzeigen mitgeteilten Elemente um $11'$ in der Rektaszension und um $2'$ in der Deklination von der Beobachtung ab. Die Ephemeride (Monatl. Corresp., Februar) gibt jene zu klein und diese zu groß an. Diese Differenzen lassen sich schon zum Teil der Wirkung der Störungen zuschreiben.

\par\indent \emph{Pallas} war schon am 20. Mai wiedergefunden worden. Schlechtes Wetter, Mondschein und ein zu großer Abstand von Sternen mit gut bestimmter Position haben aber bisher keine Beobachtungen zugelassen. Nach einer ungefähren Schätzung schien die Ephemeride im Dezemberheft der Monatl. Corresp. die Deklination gut und die Rektaszension um $2'$ zu klein anzugeben.

\par\indent \emph{Ceres} stand bisher noch zu tief in der Morgendämmerung. Wir hoffen aber, auch von ihr bald Beobachtungen mitteilen zu können.
